% ================================================================
%  chapters/ch00-intro.tex
%  导读与方法总览（不编号）
%  作用：全书的方法选择层。回答「我该用哪个方法」而不是「这个方法怎么做」。
% ================================================================

\chapterheadstar{导读与方法总览}%
{全书由 12 章方法正文、这一篇跨章导航、以及符号与索引组成。本篇给出研究问题到方法的
选择路径、12 章之间的互补关系、术语与来源约定。方法本身的原理、公式与读图逻辑在各章。}

\subsection{这本书怎么用}

本书不按原书顺序讲仪器，而按\textbf{你要回答的问题}组织。每个方法都按同一套结构展开，
所以你可以只读其中一节，也可以拿两个方法对照着读：

\begin{center}
{\srcsize
\begin{tabularx}{\columnwidth}{@{}lX@{}}
\toprule
环节 & 它回答什么 \\
\midrule
问题与输出 & 这个方法实际测量什么，能据此推断什么 \\
原理与公式 & 核心关系式及其\textbf{成立条件} \\
测量与条件 & 样品、模式、预处理、校准中哪些决定结果 \\
数据到结论 & 从原始信号到结论的处理顺序与读图步骤 \\
误判与边界 & 干扰、模型不唯一性、不能独立支持的结论 \\
互补与回查 & 该与哪个方法联用，原书精确定位在哪 \\
\bottomrule
\end{tabularx}}
\end{center}

读任何一节时，请至少读完「误判与边界」。绝大多数错误结论不是因为不会测量，
而是因为忽略了成立条件。

\subsection{五个问题}

本书的目标是让读者能回答企划约定的五个问题：

\begin{enumerate}
\item 这个方法实际测量什么，可据此推断什么？（读每节的「问题与输出」）
\item 核心原理、公式及其成立条件是什么？（「原理与公式」，条件写在公式旁）
\item 实验数据、谱图或图像应当怎样分析？（「数据到结论」）
\item 结论可能受哪些假设、干扰或处理步骤影响？（「误判与边界」）
\item 需要哪一种互补方法进一步验证？（「互补与回查」与下节的矩阵）
\end{enumerate}

\subsection{方法选择矩阵}

\begin{table}[tbp]
\centering
\caption{研究问题到方法的选择矩阵。「首选」是能独立支撑结论的方法，「互补」用于交叉验证或补足首选方法的盲区。
引用的是本书章号，不是原书章号\src{本表为本改编本归纳，非源书表格；原书对应论述见 S1 §1.1.2，书内 p.17 / PDF p.36。}}\label{tab:selector}
{\srcsize
% 【列宽修正】原为 {p{46mm}p{30mm}p{42mm}X}：前三列定宽合计 118 mm，X 列只剩余量。
% 单栏改版后 \twocolwidth 由 176 mm 收到 158 mm，X 列被压到 40 mm，多个单元格内容
% 超出其定宽而溢出（实测 p.5 出现 XPS 与 NMR 两串字符叠印）。
% 改为四个等宽 X 列，由 tabularx 自动分配，任何版心宽度下都不会溢出。
\begin{tabularx}{\twocolwidth}{@{}>{\raggedright\arraybackslash}X
                                  >{\raggedright\arraybackslash}X
                                  >{\raggedright\arraybackslash}X
                                  >{\raggedright\arraybackslash}X@{}}
\toprule
要回答的问题 & 首选方法 & 互补方法 & 为什么需要互补 \\
\midrule
体相是什么物相、晶胞参数多少 & 多晶 XRD（\chapref{ch:04}） & 中子衍射、\ce{Raman}（\chapref{ch:06}） & XRD 对轻元素与无序结构不敏感 \\
比表面积、孔容与孔径分布 & 低温氮吸附（\chapref{ch:01}） & 压汞法、电子显微统计（\chapref{ch:02}） & 吸附法给等效孔宽，显微给真实形貌 \\
表面有哪些元素、什么化学态 & XPS（\chapref{ch:09}） & AES、UPS、ISS（\chapref{ch:09}） & XPS 定量为半定量，深度不均匀会失真 \\
活性位的数量与强度 & 化学吸附 + TPD（\chapref{ch:05}） & 红外探针分子（\chapref{ch:07}）、\ce{1H} MAS NMR（\chapref{ch:08}） & 吸附法不分位点种类，探针分子可分辨 Br\o nsted/Lewis \\
局域配位与短程有序 & 固体 MAS NMR（\chapref{ch:08}） & \ce{Raman}（\chapref{ch:06}）、XAFS（本书未系统介绍） & NMR 只对特定核灵敏，需元素互补 \\
颗粒尺寸与形貌 & 透射电子显微（\chapref{ch:02}） & XRD 峰宽（\chapref{ch:04}）、STM/AFM（\chapref{ch:12}） & 显微统计量小，衍射给的是体积平均 \\
表面吸附物种与反应中间物 & 原位红外（\chapref{ch:07}） & 原位 Raman（\chapref{ch:06}）、原位 NMR（\chapref{ch:08}） & 各方法窗口不同：红外怕水、Raman 怕荧光、NMR 怕顺磁 \\
反应速率、机理与动力学参数 & 动力学测定（\chapref{ch:10}） & SSITKA、TAP（\chapref{ch:10}） & 稳态动力学对机理是间接的，需非稳态方法证中间物 \\
电催化活性与稳定性 & 电化学方法（\chapref{ch:11}） & 原位红外/Raman/质谱（\chapref{ch:11}） & 电化学只给总电流，需谱学识别中间物 \\
原子尺度表面结构与成键 & STM/AFM（\chapref{ch:12}） & TERS、Nano-FTIR（\chapref{ch:12}） & 显微给几何，近场光谱给化学识别 \\
热稳定性、积炭与失活 & TG/DTA/DSC（\chapref{ch:03}） & TPO（\chapref{ch:05}）、\ce{13C} NMR（\chapref{ch:08}） & 热分析给总量与温度，不给积炭类型 \\
\bottomrule
\end{tabularx}}
\end{table}

\begin{crosscheck}
上表的「首选」不是「最好」，而是\textbf{在该问题下最能独立支撑结论}。任何单一方法都
只能给出一个侧面的证据；凡是把「拟合较好」「单个峰归属」「多方法间相关」直接当作唯一
证明的写法，本书一律不采用。判断方法之间是否真的互补，标准是它们的\textbf{物理量不同}，
而不是仪器不同——两台都给元素组成的仪器不构成互补。
\end{crosscheck}

\subsection{选择路径}

不知道从哪开始，按图~\ref{fig:chooseroute} 的顺序问三个问题。

\begin{figure}[tbp]
\centering
\begin{tikzpicture}[
  x=1mm,y=1mm,
  box/.style={draw=accent,line width=0.5pt,rounded corners=1pt,align=center,
              inner sep=1.4mm,font=\fontsize{8}{10}\selectfont,text width=21mm},
  dec/.style={draw=accent,line width=0.5pt,align=center,inner sep=1.5mm,
              font=\fontsize{8}{10}\selectfont,text width=21mm,
              shape=diamond,aspect=2.1},
  ans/.style={align=center,font=\fontsize{8}{10}\selectfont,text width=21mm},
  ar/.style={-{Stealth[length=2mm]},line width=0.6pt,draw=accent}
]
% 【重叠修正·第二轮】原来各子节点用 below left/right=<v> and <h> of <父> 定位——
% 那定位的是**节点角点**，同一父节点的多个子节点会从父节点中心向两侧铺开而互相侵入。
% 实测 p.6 有 4 对框重叠，最严重一对 66.9x28.4 pt（几乎完全重合）。
% 现在全部改用**显式 (x,y) 坐标**：相邻框中心间距 28.8 mm
% = 框宽 23.8 mm（text width 21 + 2x1.4 内边距）+ 5 mm 缝，是算出来的确定值。
% 图幅 x 从 -67.3 到 81.7 mm = 149 mm，小于单栏版心 165 mm。
\node[dec] (q1) at (0,0) {要问的是\\体相还是表面？};
\node[box] (bulk) at (-41,-18) {体相结构\\晶相、晶胞、价态};
\node[box] (surf) at (41,-18) {表面与界面\\组成、位点、吸附态};
\node[dec] (q2) at (-41,-37) {有序还是\\无定形/短程？};
\node[dec] (q3) at (41,-37) {元素、位点\\还是吸附态？};
\node[box] (xrd) at (-55.4,-53) {多晶 XRD\\\chapref{ch:04}};
\node[box] (nmr) at (-26.6,-53) {MAS NMR\\\chapref{ch:08}\\XAFS（未系统介绍）};
\node[box] (xps) at (12.2,-53) {XPS/AES\\\chapref{ch:09}};
\node[box] (ads) at (41,-53) {化学吸附/TPD\\\chapref{ch:05}};
\node[box] (insitu) at (69.8,-53) {原位 IR/Raman\\原位 NMR\\\chapref{ch:07}\chapref{ch:06}};

\draw[ar] (q1) -- (bulk);
\draw[ar] (q1) -- (surf);
\draw[ar] (bulk) -- (q2);
\draw[ar] (q2) -- (xrd);
\draw[ar] (q2) -- (nmr);
\draw[ar] (surf) -- (q3);
\draw[ar] (q3) -- (xps);
\draw[ar] (q3) -- (ads);
\draw[ar] (q3) -- (insitu);

\node[ans] at (-55.4,-68) {峰位$\to$物相\\峰宽$\to$尺寸};
\node[ans] at (-26.6,-68) {配位\\短程有序};
\node[ans] at (12.2,-68) {元素\\化学态};
\node[ans] at (41,-68) {位点数量\\强度};
\node[ans] at (69.8,-68) {中间物\\反应路径};
\end{tikzpicture}
\caption{方法选择路径。读者应看的是\textbf{分叉的依据}：先分体相/表面，再分「有序程度」或
「要元素、位点还是吸附态」。图中只给最常用的入口，边界情况见正文各章的「互补与回查」
与表~\ref{tab:selector}。\schemamark}
\label{fig:chooseroute}
\end{figure}

\subsection{术语、符号与来源约定}

\begin{itemize}
\item \textbf{缩写}：首次出现给出中文名、英文全称与缩写，后文统一用缩写，保留英文检索入口。
完整缩写表见书末「符号与缩写」。
\item \textbf{符号}：同一符号在不同章可能含义不同（例如 $C$ 在吸附章是 BET 常数、
在动力学章可能指浓度），各章在局部重新定义，不跨章沿用。
\item \textbf{来源标记}：正文用 8\,pt 的方括号短标，形如
\src{S1 §1.1.2，书内 p.17 / PDF p.36，式(1-13)}。
其中 \textbf{S1 = 上册}，\textbf{S2 = 下册}；同时给出书内页码与 PDF 页序，
因为全书页码映射分段，不能套用固定偏移。
\item \textbf{内容性质}：凡不是原书实验结果的数值，一律标 \pedset{教学设定}；
凡原书之外的后续发展，标 \supp{补充} 并给独立来源与日期；
凡尚未对照原页面图像核对的公式，标 \tocheck{待核对}。
\item \textbf{核对状态}：来源标记只说明内容出自原书的哪一处，供回查；本版不标「已核对」。
图与表若有改编、增删列，或含本改编本自己设定的数值，一律在题注内写明。
\item \textbf{图源四类}分别标注：原书实测图（注明册别、书内页码、PDF 页序、原图号）、
\redrawmark（据原书改绘）、\modelmark（模型计算）、\schemamark（示意）。
\end{itemize}

\begin{pitfall}
不要用本书的「来源标记」当作二次文献引用。本书的来源指向原书章节与页码，
用于回查；若你要在论文中引用原始工作，请循原书参考文献找到原始文献并自行核对。
本书不做转引背书。
\end{pitfall}

\subsection{三条阅读路线}

\begin{description}[leftmargin=0pt,style=nextline,font=\heiti]
\item[快速查阅] 用表~\ref{tab:selector} 定位方法，读该章「问题与输出」「误判与边界」两节。
\item[系统学习] 按章顺序读；跨章重复的原理只在主定义处展开，其余以提醒与交叉引用带过。
\item[方案设计] 先读本篇与第 10 章（动力学）确立「要测什么量」，
再按测得量的物理性质选择表征手段，最后用「互补与回查」搭出证据链。
\end{description}
